% !TEX program = xelatex
\documentclass[cn, chinesefont=nofont, color=black, device=normal, 11pt]{elegantnote}

% 中文字体设置 (遵循 skill 规范)
\setCJKmainfont[AutoFakeSlant=0.2]{Noto Serif CJK SC}
\setCJKsansfont{Noto Sans CJK SC}
\setCJKmonofont{Noto Sans CJK SC}

\title{Vision\_Board 研发日志}
\author{Wenbo}
\date{2026-04-13}

\begin{document}
\maketitle

\section{今日目标}
\begin{itemize}
    \item 完善 LED 示例应用，验证面向对象线程管理框架的正确性。
    \item 确保主线程正确挂起，防止空闲线程栈溢出崩溃。
    \item 验证 FreeRTOS 线程调度与 RT-Thread 的共存稳定性。
\end{itemize}

\section{今日进展}
本次提交完成了 LED 示例应用的完整实现，遵循嵌入式 C++ 高阶编程规范：

\subsection{核心代码实现}
\begin{enumerate}
    \item \textbf{面向对象线程封装}（\texttt{src/app/app\_led.cpp/h}）：
          \begin{itemize}
              \item 类内部自行管理线程句柄与栈空间，外部调用者无需感知线程细节。
              \item 使用 \texttt{rt\_thread\_t} 管理线程生命周期，调用 \texttt{start()} 即自动创建并启动线程。
              \item 线程函数声明为 \texttt{private}，防止外部误操作。
          \end{itemize}
    \item \textbf{静态内存管理}：
          \begin{itemize}
              \item 全局静态实例 \texttt{AppLED sys\_led} 分配在 .bss/.data 段，不使用堆内存。
              \item 线程栈大小固定为 512字节，符合嵌入式资源受限场景。
          \end{itemize}
    \item \textbf{HAL 层入口修复}（\texttt{src/main.cpp}）：
          \begin{itemize}
              \item 修复 \texttt{hal\_entry} 函数：主线程不再盲目返回，而是进入 \texttt{while(1)} 永久挂起。
              \item 彻底根除因主线程提前退出、\texttt{tidle0} 空闲线程执行清理导致 Stack Overflow 的问题。
          \end{itemize}
\end{enumerate}

\subsection{Git 提交记录}
\begin{itemize}
    \item \texttt{494d2e7} — \texttt{feat(app): 添加 LED 示例应用}
\end{itemize}

\section{项目进度概览}
参照 \texttt{docs/Project\_Roadmap\.md}，当前各模块完成度如下：

\begin{table}[htbp]
    \centering
    \begin{tabular}{lp{6cm}}
        \toprule
        \textbf{模块}                  & \textbf{状态} \\
        \midrule
        网络中间件 (micro-ROS + TCP/HTTP) & [ ] 待完成     \\
        静态内存池设计                      & [ ] 待完成     \\
        零拷贝数据流转                      & [ ] 待完成     \\
        硬件 JPEG 编码器                  & [ ] 待完成     \\
        端侧 AI 推理部署                   & [ ] 待完成     \\
        优先级调度防反转                     & [ ] 待完成     \\
        ROS2 生态融合验证                  & [ ] 待完成     \\
        全景可视化监控                      & [ ] 待完成     \\
        \midrule
        LED 示例应用                     & [x] 已完成 已完成 \\
        FreeRTOS 基础集成                & [x] 已完成 已完成 \\
        RA 平台图形支持                    & [x] 已完成 已完成 \\
        RA8D1 BSP 基础框架               & [x] 已完成 已完成 \\
        \bottomrule
    \end{tabular}
\end{table}

\section{问题与反思}
\begin{itemize}
    \item 昨日发现的 \texttt{hal\_entry} 函数提前返回导致系统崩溃的问题已彻底修复。
    \item 后续新增应用模块时，应遵循相同的\"主线程永久挂起\"模式。
\end{itemize}

\section{明日计划}
\begin{itemize}
    \item 实现开发板端的摄像头驱动程序与 LCD 屏幕显示。
    \item 设计静态图像缓冲区的基础框架（SPSC 队列）。
\end{itemize}

\end{document}
